\begin{textsample}{POS Dim 1 – ma – Score 50.00 – ma/ma\_em\_45.txt}  \label{ex:f1_pos_069}
Itinerário de Ciências Exatas, Tecnológicas e da Terra Enfoque do itinerário : Matemática | Geografia | Sociologia | Biologia | Física | Química A união do homem com a máquina transformou a forma de produção de bens de consumo e alimentos, assim como o modo de vida e as relações sociais da sociedade.
Assim, a produção agrícola para manutenção da vida humana no planeta começou a usar os conhecimentos das \textbf{ciências} da natureza, da geografia e da matemática para buscar soluções aos fenômenos naturais que vinham comprometendo os resultados \textbf{esperados} na produção.
Neste itinerário, a área de matemática propõe o \textbf{aprofundamento} de conhecimentos que estruturam os conceitos matemáticos de forma significativa na formação do estudante e, por consequência, de um cidadão com perfil empreendedor, competente, com consciência socioeconômica, investigativa e ética e com condições de desenvolver e realizar seus projetos individuais e coletivos, comprometido com o desenvolvimento local e regional, por meio de um currículo que contemple o estudo \textbf{focado} na resolução de \textbf{problemas} e \textbf{análises} \textbf{complexas} ( \textbf{funcionais} e não lineares ), \textbf{análise} de dados estatísticos e probabilidade, geometria, programação, jogos \textbf{digitais}, sistemas dinâmicos, entre outros.
O mundo tecnológico despontou com a transformação \textbf{digital} que inseriu a inteligência artificial, painel análogo, big data, machine learning, IoT ( Internet das Coisas ), robôs autônomos colaborativos, realidade aumentada etc., implementados no setor industrial, no conhecimento e compartilhamento de informação.
Isto posiciona os países mais desenvolvidos no contexto \textbf{econômico} e social.
Nesse contexto, o itinerário \textbf{Ciências} Exatas, Tecnológicas e da Terra vem articular o conhecimento aos estudantes maranhenses do ensino \textbf{médio} para o desenvolvimento do exercício do pensamento investigativo, experimental, inovador, \textbf{crítico} e participativo e eticamente responsável para o exercício da cidadania e \textbf{progressão} dos seus estudos e no mundo do trabalho.
Essa integração se expressa na forma como o itinerário dá enfoque nos componentes curriculares matemática, física, química, \textbf{biologia}, geografia e sociologia, que articulam de forma interdisciplinar e transdisciplinar com as áreas de Ciências Humanas e Linguagens e de Ciências Humanas e Sociais Aplicadas.
Inserido no contexto deste itinerário formativo, destaca -se que, no campo produtivo das Ciências Exatas, Tecnológicas e da Terra, o \textbf{aprofundamento} da área de Matemática possibilita a ampliação, o \textbf{aprofundamento} e o enriquecimento do repertório de conhecimentos dos estudantes, expandindo, dessa forma, suas \textbf{capacidades} de ler o mundo de \textbf{maneira} \textbf{crítica} e propositiva e, mais ainda, de sua própria atuação como estudante, como protagonista e como agente de transformação da sociedade, proporcionando \textbf{introdução} a conhecimentos técnicos necessários para as carreiras afins desse itinerário formativo, projeto e produção em múltiplas áreas, no dimensionamento de componentes e elementos de artefatos em geral, propondo uma reflexão criativa.
Nesse contexto, o \textbf{aprofundamento} do estudo da física como componente básico do itinerário formativo das Ciências Exatas, Tecnológicas e da Terra permite ao estudante a \textbf{compreensão} dos fenômenos da natureza e a investigação de suas causas e efeitos na vida cotidiana da sociedade, de forma a levá -lo à reflexão e \textbf{análise} \textbf{crítica} dos mesmos e buscar soluções tecnológicas \textbf{digitais} de informação e comunicação, de forma a preservar a vida em todas as suas formas e os processos produtivos de \textbf{maneira} sustentável, utilizando \textbf{métodos} experimentais e \textbf{análise} de dados, com aplicações na engenharia, computação e robótica.
De forma similar, o componente curricular química busca o \textbf{aprofundamento} dos assuntos que tratam dos aspectos de exploração da terra, como estudo dos compostos químicos, preparo de soluções, ciclos biogeoquímicos, controle do pH do solo e as condições necessárias para conservação do solo para o plantio de lavouras e produtividade de organismos vivos, além da pesquisa e investigação do melhoramento dos alimentos produzidos Neste itinerário, a \textbf{biologia} é o alicerce, na medida em que reconhece a necessidade de serem fortalecidas as relações do homem com os recursos naturais e as diferentes formas de vida.
Além disso, a investigação científica é \textbf{inerente} à \textbf{biologia}, que, por meio de ensaios práticos-experimentais, alicerçada no \textbf{método} científico, dará a \textbf{compreensão} dos fenômenos ocorridos pela intervenção do homem no processo do mundo natural.
A importância da \textbf{filosofia} no estudo das Ciências Exatas, Tecnológicas e da Terra ocorre \textbf{devido} à integração e à \textbf{contextualização} histórico-social do próprio conhecimento científico.
Nesta lógica, segundo Matthews, “ a Filosofia da Ciência está vazia sem História da Ciência ; a História da Ciência está cega sem Filosofia da Ciência ” ( 1995, p. 174 ).
As descobertas da \textbf{ciência} são de grande relevância para os \textbf{debates} filosóficos sobre as diversas concepções de \textbf{ciência}, a relação mente/corpo, o \textbf{método}, os tipos, formas, origens e essência do próprio saber em geral.
Por sua vez, também é papel da \textbf{filosofia} refletir sobre os padrões de racionalidade científica, o valor cognitivo das \textbf{teorias}, os esquemas de explicação, a evolução da \textbf{ciência} e outros \textbf{problemas} semelhantes.
Ainda dentro deste contexto, o ensino de \textbf{ciência} por meio de uma \textbf{compreensão} \textbf{reflexiva} filosófica nos propõe ultrapassarmos as distorções sobre a \textbf{ciência}, aprofundar a investigação sobre a ação das \textbf{ciências} da natureza na história da \textbf{humanidade}, com a intenção de desenvolver, compreender e \textbf{aprimorar} amplamente o pensamento científico e o próprio pensar filosófico como um todo.
Junto a isso, destaca -se o conhecimento científico da geografia, que \textbf{mobiliza} a \textbf{compreensão} do espaço multidimensional ( físico e humano ).
Entender a dimensão física do espaço geográfico e seus componentes estruturais possibilitará ao estudante o desenvolvimento de habilidades voltadas à exploração racional dos recursos.
Os conhecimentos \textbf{mobilizados} no âmbito da geografia e física ( geologia, geomorfologia, pedologia, climatologia e hidrologia ) são eixos estruturantes na construção de diferentes áreas de atuação profissional nas Ciências Exatas, da Terra e Tecnológicas, como a engenharia civil, que \textbf{demandará} de conhecimentos de geologia estrutural, ou a engenharia agrária, que \textbf{exigirá} conhecimentos de pedologia e climatologia. \textbf{Verifica} -se, ainda, uma aplicação multiprofissional dos conhecimentos da cartografia, que são trabalhados na geografia para a \textbf{compreensão} de fenômenos espaciais.
Por \textbf{conseguinte}, os pressupostos \textbf{cartográficos} encontram na matemática suas bases \textbf{teóricas} e práticas, como no estabelecimento de \textbf{projeções} e escalas.
Ressalta -se, também, que a \textbf{análise} em multiescala feita pela geografia proporciona uma visão holística dos fenômenos, que, por sua vez, podem \textbf{auxiliar} na tomada de decisões em diferentes áreas.
Outra aplicação é o tratamento das informações geográficas a partir das bases fundamentais da matemática, por exemplo, nos estudos demográficos, em que o estudante realizará a imersão dos conteúdos de população associados à estatística.
Por outro \textbf{lado}, entender a organização da sociedade ( economia, política, cultura etc. ) e como ela se estabelece em seu território permitirá que o estudante problematize fenômenos sociais e construa reflexões sobre eles, como aqueles voltados ao espaço agrário brasileiro.
Em história, um aspecto marcante na matemática foi a forma geométrica das pirâmides e toda a engenharia envolvida nas suas construções, na \textbf{análise} estatística da população e das comunidades tradicionais desde o início da colonização portuguesa no Brasil.
Nisso destaca -se, também, a importância da matemática no desenvolvimento das civilizações : leitura de gráficos estatísticos para diagnosticar situações simuladas e reais e propor intervenções concretas para soluções de \textbf{problemas} sociais e \textbf{econômicos} ; o Renascimento Científico e o estudo da trigonometria como forma de entender os astros ; desigualdade socioeconômica e distribuição de renda para as populações negras, índices de desemprego e escolarização.
A sociologia com o teor do itinerário formativo de Ciências Exatas e Tecnológicas e da Terra possui uma relação direta com a produção do conhecimento sociológico, \textbf{tanto} no que se refere às questões \textbf{metodológicas} quanto à própria materialidade das relações sociais construídas em torno desse campo produtivo, que, por si só, já é o objeto de estudo dessa \textbf{ciência} ( as relações e estruturas sociais, a cultura e os sistemas políticos ).
Entende -se, pois, que a linguagem da matemática \textbf{subsidia} a construção desse pensamento científico acerca da realidade social por meio das pesquisas quantitativas.
Esse embasamento será o arcabouço desse itinerário e, assim, conduzirá os estudantes ao exercício da investigação, à confirmação das evidências científicas e à reflexão sobre as questões abordadas, \textbf{auxiliando} -o no pensamento \textbf{crítico} e a um fazer participativo e eticamente responsável.
A segunda relação se desenvolve no eixo cultura e trabalho, desdobrando -se em vários objetos de estudo, em torno da relação do ser humano com o meio ambiente, o trabalho e a cultura, pois, na medida em que transforma a natureza, ele produz cultura e os diversos modos de vida.
A temática sociedade e meio ambiente é fundamental para essa \textbf{compreensão}, pois a história das sociedades também é a história de múltiplas relações com o meio ambiente, uma vez que cada sociedade encontra uma forma específica de satisfazer suas necessidades, e estas são socialmente construídas.
Nessa perspectiva, poderão \textbf{emergir} neste itinerário o \textbf{aprofundamento} das seguintes discussões : 1. alimentação como direito humano ; 2. os impactos socioeconômicos do uso de agrotóxicos ; 3. soberania alimentar ; 4. consumo consciente para o equilíbrio ambiental ; 5. fome vs. \textbf{tecnologias}, crise alimentar e globalização.
Nesse itinerário formativo, portanto, o \textbf{aprofundamento} da área de Matemática proporcionará a \textbf{introdução} de conhecimentos técnicos necessários para as carreiras afins do campo produtivo referente ao formativo, projeto e produção em múltiplas áreas, no dimensionamento de componentes e elementos de artefatos, em geral, propondo uma reflexão criativa.
A área de Linguagens contribuirá porque a necessidade de comunicação é própria do ser humano e este, pertencendo a um contexto sócio-histórico-cultural, tem latente a necessidade de ampliar sua cosmovisão e possibilidades comunicativas, para o que recorre à aprendizagem de outros idiomas, sobretudo com a intensificação do processo de globalização.
Nesse viés, destacamos que é \textbf{imprescindível} que o ensino-aprendizagem da língua estrangeira garanta as bases \textbf{discursivas}, que é o domínio da \textbf{capacidade} de se comunicar com outras pessoas por meio do texto escrito ou oral.
Dessa forma, \textbf{tanto} o \textbf{inglês} instrumental como o \textbf{espanhol} instrumental poderão \textbf{somar} para este itinerário formativo.
O componente curricular educação física poderá proporcionar um \textbf{aprofundamento} na sua formação, contribuindo ao desenvolvimento de suas habilidades e \textbf{competências}, buscando atender às necessidades, mantendo um diálogo mais dinâmico com os estudantes por meio do estudo e da experimentação das práticas corporais.
Quanto à contribuição da arte no itinerário formativo, cabe evidenciar que a relevância deste componente curricular consiste em potencializar a \textbf{compreensão} dos objetos do conhecimento nas áreas das Ciências Exatas, Tecnológicas e da Terra – na perspectiva de potencializar a formação do educando nos diversos cursos de abrangência deste itinerário –, que se dá por meio da estesia, da catarse e da fruição em arte, próprios do processo ensino-aprendizagem em arte, conforme orienta a proposta triangular de Ana Mae Barbosa, que consiste em três pilares : o fazer, o apreciar e o contextualizar.
O componente curricular língua portuguesa, no itinerário de Ciências Exatas, Tecnológicas e da Terra, oferece a contribuição no trabalho, sobretudo com os gêneros que emanam do contexto ( semiose ) \textbf{digital}, haja \textbf{vista} o trabalho que os diferentes componentes/cursos desse itinerário \textbf{demandam} sob essa perspectiva.
Pensar dessa forma é estar alinhado com o contexto da modernidade tardia, que busca \textbf{competências} e habilidades para as próximas gerações que são intermediadas pelo universo tecnológico/digital.

% matched lemmas: análise, aprimorar, aprofundamento, auxiliar, biologia, capacidade, cartográfico, ciência, competência, complexo, compreensão, conseguinte, contextualização, crítico, debate, demandar, devido, digital, discursivo, econômico, emergir, espanhol, esperar, exigir, filosofia, focar, funcional, humanidade, imprescindível, inerente, inglês, introdução, lado, maneira, metodológico, mobilizar, médio, método, problema, progressão, projeção, reflexivo, somar, subsidiar, tanto, tecnologia, teoria, teórico, verificar, vista
\end{textsample}
